% 根据《物理实验报告模板.docx》整理的 LaTeX 源码。
% 中文排版使用 XeLaTeX。
% 请将 tongji_logo.jpeg 与本文件放在同一目录。
\documentclass[UTF8,a4paper,10pt,fontset=fandol]{ctexart}
\usepackage[left=3.175cm,right=3.175cm,top=2.54cm,bottom=2.54cm]{geometry}
\usepackage{graphicx}
\usepackage{tikz}
\usepackage{fancyhdr}
\usepackage{lastpage}

\setlength{\parindent}{0pt}
\setlength{\parskip}{0pt}
\setlength{\headheight}{18pt}
\linespread{1}
\pagestyle{empty}

% 封面填写横线。
\newcommand{\coverfield}[1]{%
  \noindent{\heiti\fontsize{18pt}{48pt}\selectfont
  \makebox[3.3cm][l]{#1}\rule[-0.25em]{8.8cm}{0.45pt}\par}%
}

% 装订线依照 Word 模板的 39 个逐行字符排布。
% 第 14、20、26 个字符分别为“装”“订”“线”，字间各有五个“┊”。
\newcommand{\bindingline}[1][2.85]{%
  \begin{tikzpicture}[remember picture,overlay]
    \foreach \i in {0,...,38}{%
      \pgfmathsetmacro{\bindingy}{-#1-0.52*\i}%
      \ifnum\i=13\relax
        \node[inner sep=0pt] at
          ([xshift=1.05cm,yshift=\bindingy cm]current page.north west)
          {\songti\fontsize{10.5pt}{10.5pt}\selectfont 装};
      \else\ifnum\i=19\relax
        \node[inner sep=0pt] at
          ([xshift=1.05cm,yshift=\bindingy cm]current page.north west)
          {\songti\fontsize{10.5pt}{10.5pt}\selectfont 订};
      \else\ifnum\i=25\relax
        \node[inner sep=0pt] at
          ([xshift=1.05cm,yshift=\bindingy cm]current page.north west)
          {\songti\fontsize{10.5pt}{10.5pt}\selectfont 线};
      \else
        \pgfmathsetmacro{\bindingtop}{\bindingy+0.20}%
        \pgfmathsetmacro{\bindingbottom}{\bindingy-0.20}%
        \draw[line width=0.35pt,dash pattern=on 1.2pt off 1.8pt]
          ([xshift=1.05cm,yshift=\bindingtop cm]current page.north west)
          -- ([xshift=1.05cm,yshift=\bindingbottom cm]current page.north west);
      \fi\fi\fi
    }
  \end{tikzpicture}%
}

\fancypagestyle{report}{%
  \fancyhf{}
  \fancyhead[C]{}
  \fancyfoot[C]{第\thepage 页\quad 共\pageref{LastPage}页}
  \renewcommand{\headrulewidth}{0.4pt}
  \renewcommand{\footrulewidth}{0pt}
}
\fancypagestyle{record}{%
  \fancyhf{}
  \fancyhead[C]{原始数据记录单}
  \fancyfoot[C]{第\thepage 页\quad 共\pageref{LastPage}页}
  \renewcommand{\headrulewidth}{0.4pt}
  \renewcommand{\footrulewidth}{0pt}
}

\begin{document}
\fontsize{10.5pt}{20pt}\selectfont

% 第 1 页：封面
\thispagestyle{empty}
\vspace*{1.0cm}
\begin{center}
  \includegraphics[width=8.184cm]{tongji_logo.jpeg}\par
  {\bfseries\fontsize{18pt}{24pt}\selectfont TONGJI UNIVERSITY\par}
  \vspace{1.5cm}
  {\heiti\fontsize{36pt}{44pt}\selectfont 物理实验报告\par}
\end{center}

\vspace{1.5cm}
\coverfield{实验名称}
\coverfield{学生姓名}
\coverfield{学\hspace{1em}号}
\coverfield{组\hspace{1em}号}
\coverfield{实验日期}

\clearpage

% 第 2 页：实验报告正文
\thispagestyle{report}
\bindingline
\begin{center}
  {\heiti\fontsize{14pt}{40pt}\selectfont 实验名称\par}
\end{center}

\begingroup
\setlength{\parindent}{0pt}
\fontsize{10.5pt}{20pt}\selectfont
{\heiti 实验目的}\par
\vspace{3\baselineskip}

{\heiti 实验原理}\par
（请对实验原理做简要概述，需要包含必要的原理图、公式等，不超过300字。）\par
\vspace{2\baselineskip}

{\heiti 实验内容}\par
\vspace{3\baselineskip}

{\heiti 实验仪器}\par
\vspace{3\baselineskip}

{\heiti 数据记录处理}\par
\vspace{3\baselineskip}

{\heiti 讨论与分析}\par
\vspace{4\baselineskip}

{\heiti 实验结论}\par
\endgroup

\clearpage

% 第 3 页：原始数据记录单
\thispagestyle{record}
\bindingline[3.03]
\vspace*{0.6cm}
（请将教师签字的实验数据记录单完整拍照粘贴至此处。）

% 如需插入照片，可把下行取消注释，并将文件名替换为实际照片：
% \begin{center}\includegraphics[width=\textwidth,height=19cm,keepaspectratio]{原始数据记录单.jpg}\end{center}

\end{document}
